%! AP Statistics — Unit 4, Lesson 4.6 — Homework KEY
\documentclass[10pt]{article}
\usepackage{apstats-article}
\usepackage{apstats-key}

\begin{document}

\pageheader{Unit 4, Lesson 4.6}{Homework --- KEY}
\namedateperiod

\vspace{0.1in}

\begin{practicebox}
\small
Streaming service $\times$ age group (300 subscribers):

\vspace{4pt}
\renewcommand{\arraystretch}{1.4}
\begin{tabularx}{\linewidth}{@{} l X X X r @{}}
 & \textbf{Movies} & \textbf{TV Series} & \textbf{Documentaries} & \textbf{Total} \\
\hline
\textbf{Under 25}  & 50 & 60 & 10 & 120 \\
\textbf{25--44}    & 40 & 50 & 30 & 120 \\
\textbf{45+}       & 20 & 20 & 20 & 60  \\
\hline
\textbf{Total}     & 110 & 130 & 60 & 300 \\
\end{tabularx}

\vspace{8pt}
\textbf{1.}
\begin{enumerate}[label=(\alph*), leftmargin=2em, itemsep=4pt]
  \item \ans{$P(\text{Movies}) = 110/300 \approx 0.367$}
  \item \ans{$P(\text{Movies}|\text{Under 25}) = 50/120 \approx 0.417$}
  \item \ans{Since $0.367 \neq 0.417$, Movies and Under 25 are \textbf{not independent} --- subscribers under 25 are more likely to watch movies than the overall subscriber population.}
\end{enumerate}

\vspace{6pt}
\textbf{2.} \ans{$P(\text{TV Series}) = 130/300$; $P(\text{25--44}) = 120/300$; $P(\text{TV Series} \cap \text{25--44}) = 50/300$.\\
$P(\text{TV Series} \cup \text{25--44}) = 130/300 + 120/300 - 50/300 = 200/300 \approx 0.667$}

\vspace{6pt}
\textbf{3.}
\begin{enumerate}[label=(\alph*), leftmargin=2em, itemsep=4pt]
  \item \ans{Yes, mutually exclusive. Each subscriber has exactly one primary content type; no subscriber is counted in both Movies and Documentaries.}
  \item \ans{$P(\text{Movies} \cup \text{Docs}) = 110/300 + 60/300 - 0 = 170/300 \approx 0.567$}
\end{enumerate}

\vspace{6pt}
\textbf{4.}
\begin{enumerate}[label=(\alph*), leftmargin=2em, itemsep=4pt]
  \item \ans{No, the reasoning is incorrect. Independence is not about causation; it is about whether knowing one event changes the probability of another: the test is $P(A|B) = P(A)$.}
  \item \ans{$P(\text{TV Series}) = 130/300 \approx 0.433$; $P(\text{TV Series}|\text{Under 25}) = 60/120 = 0.5$. Since $0.433 \neq 0.5$, TV Series and Under 25 are \textbf{not independent}.}
\end{enumerate}

\vspace{6pt}
\textbf{5.} \ans{$P(\text{Docs}) = 60/300 = 0.2$; $P(\text{Docs}|\text{45+}) = 20/60 \approx 0.333$.\\
Since $0.2 \neq 0.333$, Documentaries and 45+ are \textbf{not independent}. Subscribers aged 45 and older watch documentaries at a higher rate (33.3\%) than the overall subscriber base (20\%), indicating that age and content-type preference are related for this group.}

\vspace{6pt}
\textbf{6.} \ans{True. If $P(A|B) = P(A)$, then by the definition of conditional probability: $P(A|B) = P(A \cap B)/P(B)$, so $P(A \cap B) = P(A) \cdot P(B)$. Then $P(B|A) = P(A \cap B)/P(A) = P(A) \cdot P(B)/P(A) = P(B)$. So $P(B|A) = P(B)$. Independence is a symmetric relationship: if $A$ is independent of $B$, then $B$ is independent of $A$.}
\end{practicebox}

\end{document}
